\input{../tex-inputs/latex-leseansicht-vorspann}

\section[Arthur Schnitzler an Gustav Schwarzkopf, {[}6. 4. 1903?{]}]{L04198 Arthur Schnitzler an Gustav Schwarzkopf, {[}6. 4. 1903?{]}}
\nopagebreak\mylabel{L04198v}
\rehead{ }\normalsize\beginnumbering\briefempfaengerindex{Schwarzkopf, Gustav@\textsc{Schwarzkopf, Gustav}!zzzSchnitzler, Arthur@\emph{von Arthur Schnitzler}!1903-04-062@{{[}6. 4. 1903?{]}}|(be}
\toendnotes[C]{\smallbreak\pagebreak[2]}
\correspDesc{Versand  durch Arthur Schnitzler am [6. 4. 1903?] in Edmund-Weiß-Gasse 7
\newline{}Erhalt  durch Gustav Schwarzkopf im Zeitraum [7. 4. 1903 – 11. 4. 1903?] in Tiefer Graben 23}\toendnotes[C]{\smallbreak}
\Standort{CUL, Schnitzler, B 96.}
\physDesc{Brief, 1 Blatt, 2 Seiten, 233 Zeichen
\newline{}Handschrift: Bleistift, deutsche Kurrent}\toendnotes[C]{\smallbreak}
\pstart
           \raggedleft{}{\pb}\uline{Montag.}\pend
           \vspace{0.5em}
\pstart
           lieber Guſtav, warum haben Sie nicht ein bischen gewartet! oder ſind
               wieder gekommen?\pend
           
\pstart
           Wir fahren Mittwoch; wollen Sie nicht \label{K_L04198-1v}\edtext{Dinſtag{ }Abend, (7.}{\lemma{\textnormal{\emph{Dinstag Abend, (7.}}}\Cendnote{\textnormal{Die vorliegende Karte ist, bis auf die Angabe des implizite Angabe des Wochentags (»Montag, der 6.«) undatiert. Ein Abgleich mit den in Frage kommenden
               Monaten ergibt nur den 6. 4. 1903 als möglichen Termin, der kurz vor einer Abreise stand. Auch kam es 
                  am 7. 4. 1903 zu einem Treffen mit Schwarzkopf\pwindex{Schwarzkopf, Gustav 7.\,11.\,1853 Wien – 13.\,11.\,1939 ebd.@\textsc{Schwarzkopf, Gustav} (7.\,11.\,1853 Wien – 13.\,11.\,1939 ebd.)|pwk}, das aber 
                  bereits am Nachmittag stattgefunden haben dürfte. Laut \emph{Tagebuch}\pwindex{Schnitzler, Arthur 15. 5. 1862 Wien – 21. 10. 1931 ebd.@\textsc{Schnitzler, Arthur} (15. 5. 1862 Wien – 21. 10. 1931 ebd.)!Tagebuch@\strich\emph{Tagebuch}|pwk} kam der Entschluss der Reise 
                  erst an diesem Tag zu Stande, was der Datierung widerspricht und fand die Reise ohne Olga\pwindex{Schnitzler, Olga 17.\,1.\,1882 Wien – 13.\,1.\,1970 Lugano@\textsc{Schnitzler, Olga} (17.\,1.\,1882 Wien – 13.\,1.\,1970 Lugano)|pwk} statt. Für die Datierung
                  eingewandt 
                  werden kann, dass Schnitzler am XXXX Auszeichnungsfehler: Dokument L01287 nicht gefunden Reisepläne äußerte.}}}\label{K_L04198-1}) noch ein{\pb}mal bei uns sein?\pend
           
\pstart
           Antwort überflüſſig, de{\geminationn} hoffentlich ko{\geminationm}en Sie.\pend
           
\pstart
           Herzlichſt{\\[\baselineskip]} Ihr{\\[\baselineskip]}\spacefill\mbox{A.}\pend
           \leftskip=0em{}\selectlanguage{ngerman}\endnumbering\briefempfaengerindex{Schwarzkopf, Gustav@\textsc{Schwarzkopf, Gustav}!zzzSchnitzler, Arthur@\emph{von Arthur Schnitzler}!1903-04-062@{{[}6. 4. 1903?{]}}|)be}\mylabel{L04198h}
\begin{anhang}
\end{anhang}\newcommand{\dateiname}{L04198}\newcommand{\titel}{Arthur Schnitzler an Gustav Schwarzkopf, [6. 4. 1903?]}\newcommand{\editorInnen}{Herausgegeben von Jahnke, SelmaMüller, Martin Anton}\input{../tex-inputs/latex-leseansicht-abspann}
